% Springer LNCS format - compile with pdflatex + bibtex (splncs04.bst)
\documentclass[runningheads]{llncs}

\usepackage[T1]{fontenc}
\usepackage{graphicx}
\usepackage{booktabs}
\usepackage{array}
\usepackage{amsmath}
\usepackage{hyperref}
\usepackage{url}
\renewcommand\UrlFont{\color{blue}\rmfamily}

\begin{document}

\title{CertiVerify: An Agentic Multimodal Framework for Explainable Verification of Student Activity Certificates in Indian Higher Education}

\titlerunning{CertiVerify}

\author{Leelambika K V \and Aman Singh\orcidID{} \and Sakshi P C \and Kushi S K}
\authorrunning{L. K V et al.}

\institute{Presidency University, Bangalore, India \\
\email{\{amansingh160704, pcsakshi04, kushigowdas324\}@gmail.com}}

\maketitle

\begin{abstract}
The verification of student co-curricular certificates in Indian higher education is a manual, error-prone process that undermines the integrity of records used for NAAC and NIRF accreditation. Rule-based and OCR-only pipelines fail to detect subtle forgeries, cross-field semantic mismatches, and duplicate submissions at scale. We present \textbf{CertiVerify}, an agentic multimodal framework that cross-validates certificates across five independent dimensions---text semantics, visual forensics, document layout, file metadata, and field-level cross-checking against student-submitted form data---coordinated by a central Reasoning Agent that produces a calibrated confidence score and a natural-language explanation. We also introduce \textbf{CertiForge-HE}, a benchmark of 5{,}800 Indian academic certificates across eight activity categories and four authenticity classes. Under 5-fold stratified cross-validation, CertiVerify achieves \textbf{94.1\,\% ($\pm$0.8\,\%) accuracy} and \textbf{0.95 AUC-ROC}, outperforming a fine-tuned ViT-Base baseline by 7.8 points and zero-shot multimodal LLMs by 18--24 points, while reducing estimated faculty review workload by 71\,\%.

\keywords{Certificate Verification \and Agentic AI \and Multimodal Forensics \and NER \and Explainable AI \and Indian Higher Education.}
\end{abstract}

\section{Introduction}
\label{sec:intro}

Indian higher education---over 1{,}000 universities and 40{,}000 colleges under NAAC, NIRF, and AICTE---requires institutions to document student co-curricular engagement with verifiable evidence. In practice, students upload certificates and faculty verify them manually: a slow, inconsistent process that is entirely unequipped to detect modern document manipulation. With accessible image editors and diffusion-based inpainting tools, fabricating or altering a convincing certificate requires no specialist skill.

Indian academic certificates are uniquely challenging relative to prior document-forensics work on Western invoices or ID documents: templates vary across tens of thousands of issuers; content mixes English with Devanagari, Kannada, Tamil, Telugu, and Bengali scripts; submissions arrive as smartphone photos, scans, screenshots, or PDFs with inconsistent quality; Indian names exhibit transliteration ambiguity (``Aman Singh'' vs.\ ``Amann Singh'' vs.\ ``\textit{aman sinh}''); and issuer trust varies from near-machine-verifiable (NPTEL) to untraceable (one-day workshops). Off-the-shelf forgery detectors trained on standardised Western corpora transfer poorly to this domain.

\subsubsection{Contributions.}
(1) \textbf{CertiVerify}, an agentic multimodal verification framework combining deterministic forensics, supervised classification, and learned ensemble reasoning for Indian academic certificates. (2) \textbf{CertiForge-HE}, a dataset of 5{,}800 annotated certificates across eight categories and four authenticity classes. (3) A domain-specific NER model with 0.92 macro F1 on five certificate entity types. (4) Explainable verdicts with per-agent evidence suitable for institutional audit. (5) A holistic Student Engagement Score (SES) grounded in verified records.

\section{Related Work}
\label{sec:related}

Prior student information and e-portfolio systems~\cite{balasubramanian2017,khandare2021,alqahtani2020} centralise records but rely on self-reported, unverified data. Classical document forensics uses Error Level Analysis and PDF structural checks; these fail against diffusion-based inpainting whose JPEG residuals are near-indistinguishable from genuine artifacts. TruFor~\cite{guillaro2023} applies noiseprint CNN features to splicing detection; CAT-Net~\cite{kwon2022} detects double-JPEG compression at the DCT level; ForgeryGPT~\cite{zhang2024} and M2F2-Det~\cite{guo2025} apply multimodal LLMs to forgery localisation. None are trained on academic certificates, which exhibit distinctive visual vocabulary (institutional watermarks, signature blocks, formatted date fields). Zero-shot multimodal LLM benchmarks~\cite{chen2025} report AUC 0.6--0.8 on transactional documents but cannot be deployed offline---a binding constraint under the Indian DPDP Act 2023~\cite{dpdp2023}. Recent agentic forensics decomposes document analysis across specialist modules~\cite{docagent2025}; CertiVerify adapts this pattern with two domain-specific contributions absent in prior work: a Field Cross-Check Agent that compares certificate content against student submission metadata, and calibrated trust-tier weighting based on issuer characteristics.

\begin{table}[t]
\caption{Coverage gaps in prior work.}
\label{tab:gaps}
\centering
\begin{tabular}{lll}
\toprule
\textbf{Approach} & \textbf{Coverage} & \textbf{Limitation} \\
\midrule
Activity tracking systems & Record logging & No certificate verification \\
ELA + metadata forensics & Visual tampering & Fails on AI inpainting \\
TruFor, CAT-Net & Image splicing & Not trained on certificates \\
Zero-shot MLLMs & Multimodal reasoning & No privacy, hallucination \\
Blockchain verification & Tamper-proof anchoring & Requires issuer participation \\
Generic agentic forensics & General documents & No field-level cross-check \\
\bottomrule
\end{tabular}
\end{table}

\section{The CertiVerify Framework}
\label{sec:framework}

CertiVerify decomposes verification across seven agents (Fig.~\ref{fig:arch}): one Ingestion Agent, five specialist analysers executed in parallel, and a central Reasoning Agent. All agents share a common interface---consuming a shared state object and emitting a structured evidence report with a scalar confidence in $[0,1]$ and a list of named evidence items used by the Reasoning Agent for natural-language explanation.

\begin{figure}[t]
\centering
\begin{verbatim}
    Upload -> Ingestion (OCR, Layout, Metadata, DPI)
                 |
         +-------+-------+-------+--------+
         v       v       v       v        v
       Text  Visual   Layout  Metadata  Field
                              Cross-Check
         +-------+-------+-------+--------+
                         |
                    Reasoning Agent
                 (calibrated ensemble
                  + NL explanation)
                         |
            auto_verified / needs_review / rejected
\end{verbatim}
\caption{CertiVerify multi-agent architecture.}
\label{fig:arch}
\end{figure}

\subsubsection{Ingestion Agent.} PDFs are rasterised at 300~DPI (PyMuPDF). EXIF metadata is extracted from images before any processing that might strip it. PaddleOCR~\cite{du2020} produces multi-script text (English, Devanagari, Kannada, Tamil, Telugu) and a layout map with per-region bounding boxes, font-size estimates, and OCR confidence.

\subsubsection{Text Agent.} A custom NER model (Sect.~\ref{sec:ner}) extracts five entity types. A fine-tuned sentence classifier scores semantic coherence, flagging implausible combinations (e.g., event descriptions inconsistent with declared category). A RoBERTa classifier trained on genuine and LLM-generated certificate text assigns a machine-generation probability.

\subsubsection{Visual Forensics Agent.} Error Level Analysis (re-save at JPEG Q85, pixel-wise residual) is mapped to per-region tampering probabilities by a fine-tuned ResNet-50~\cite{he2016}. Noise inconsistency analysis flags statistically mismatched regions. Institutional logos are matched against a curated database of 340 Indian HE logos.

\subsubsection{Layout Agent.} Font-size and family variance are computed across text blocks. A fine-tuned LayoutLMv3~\cite{xu2022} scores category-specific template conformance against learned spatial signatures for each of the eight activity categories.

\subsubsection{Metadata Agent.} PDF metadata (author, creator, creation/modification dates) is checked for internal consistency. EXIF fields flag editor signatures. DCT coefficient histogram analysis detects double-JPEG compression---a signature of screenshots re-edited and re-saved.

\subsubsection{Field Cross-Check Agent.} This agent---unique to certificate verification---compares extracted certificate entities against student-submitted form fields. Name pairs use Jaro-Winkler with transliteration normalisation; dates use normalised parsing with $\pm 3$ day tolerance; institutions use token overlap plus fuzzy matching; event titles use TF-IDF cosine similarity; categories use rule-based mapping. A Random Forest (500 trees, trained on 3{,}200 faculty-labelled pairs) combines per-field similarities into an overall match score.

A two-stage duplicate detector addresses a known failure mode of naive vector similarity on highly standardised documents (e.g., two certificates from the same multi-day workshop differing only in date). \textit{Stage 1:} pgvector cosine similarity $> 0.92$ on NER entity embeddings flags candidates. \textit{Stage 2:} for each candidate, exact comparison of DATE and EVENT entities---if date distance exceeds 7 days or event token overlap is below 0.6, the duplicate penalty is bypassed. Only when both stages confirm a match is \texttt{duplicate\_flag} set.

\subsubsection{Reasoning Agent.} The final score is a weighted ensemble:
\begin{equation}
s = 0.25\,t + 0.25\,v + 0.20\,\ell + 0.15(1 - m) + 0.15\,f
\label{eq:ensemble}
\end{equation}
where $t, v, \ell, m, f$ are the Text, Visual, Layout, Metadata anomaly, and Field-match confidences. Weights were determined by Bayesian optimisation on the validation fold. If \texttt{duplicate\_flag} is set, $s$ is reduced by 0.30 and routed to \texttt{needs\_review}. Verdicts: $s \geq 0.85$ $\rightarrow$ \texttt{auto\_verified}; $0.55 \leq s < 0.85$ $\rightarrow$ \texttt{needs\_review}; $s < 0.55$ $\rightarrow$ \texttt{rejected}. A structured natural-language explanation cites per-agent evidence (e.g., ``Visual Agent flagged elevated ELA residuals at the date field; Field Cross-Check confirmed name match but date mismatch of 12 days'') and is persisted in the audit log.

\section{Domain-Specific NER Model}
\label{sec:ner}

The NER corpus comprises 4{,}200 annotated certificate images: consented genuine samples across three academic years, template-generated synthetic samples, and controlled text-perturbation augmentations. Four annotators labelled under a five-entity schema (\texttt{STUDENT\_NAME}, \texttt{INSTITUTION}, \texttt{EVENT}, \texttt{DATE}, \texttt{CERT\_TYPE}) using Prodigy; inter-annotator agreement (Cohen's $\kappa$) was 0.91. The base model is spaCy \texttt{en\_core\_web\_trf}~\cite{spacy2017}, fine-tuned for 30 epochs (lr 2e-5, batch 16, patience 5). A token-level script classifier routes Devanagari and Dravidian-script tokens through IndicBERT-base~\cite{indic2020} before fusion with the spaCy transformer output. Macro F1 under 5-fold CV is \textbf{0.92}, with per-entity F1 ranging from 0.87 (EVENT) to 0.96 (DATE).

\section{CertiForge-HE Benchmark}
\label{sec:dataset}

No public Indian academic certificate dataset exists. CertiForge-HE comprises 5{,}800 samples across eight activity categories and four authenticity classes (Table~\ref{tab:dataset}).

\begin{table}[t]
\caption{CertiForge-HE composition (stratified 70/15/15 split).}
\label{tab:dataset}
\centering
\begin{tabular}{lrrrr}
\toprule
\textbf{Class} & \textbf{Train} & \textbf{Val} & \textbf{Test} & \textbf{Total} \\
\midrule
Genuine     & 2{,}030 & 435 & 435 & 2{,}900 \\
Forged      &    820 & 175 & 175 & 1{,}170 \\
AI-Edited   &    710 & 152 & 152 & 1{,}014 \\
Duplicate   &    500 & 107 & 107 &    714 \\
\bottomrule
\end{tabular}
\end{table}

\textit{Genuine} samples were student-contributed (2022--2024, Presidency University) with written informed consent. \textit{Forged} samples were produced via GIMP, Photoshop, and Canva using logged manipulation types (text replacement, logo substitution, border modification). \textit{AI-Edited} samples used Stable Diffusion XL inpainting for text-region replacement and Gemini 2.5 Flash for full generation. \textit{Duplicate} samples include rotation, compression, and crop variants.

\subsubsection{Ethics and Consent.} The study was approved by the Department of Computer Science Ethics Committee, Presidency University. PII in released samples is pseudonymised; small-organiser logos are blurred unless consent was obtained. The dataset will be released under a research-only licence requiring institutional sign-off.

\section{Deployment}
\label{sec:deployment}

Training used a single NVIDIA A100 (40~GB). Inference runs on two NVIDIA T4 GPUs (16~GB), with ONNX Runtime INT8 quantisation reducing per-model memory by $\approx$55\,\% at accuracy cost below 0.4 percentage points. Orchestration uses a LangGraph cyclic state graph: the Ingestion Agent populates shared state, the five specialists run as concurrent nodes, and the Reasoning Agent is gated on all five \texttt{DONE} signals. Wall-clock latency is bounded by the slowest agent (Visual Forensics, $\approx 2.1$~s); end-to-end $\approx 2.6$~s per certificate. Under 50 simultaneous uploads, p95 remains at 2.9~s and p99 at 3.4~s. The web stack is React 18 + TypeScript with Supabase (PostgreSQL + Auth + Storage) enforcing row-level security.

\section{Results}
\label{sec:results}

All results are 5-fold stratified cross-validation on CertiForge-HE train+validation, with final evaluation on the held-out test fold ($n = 869$).

\begin{table}[t]
\caption{End-to-end verification performance.}
\label{tab:main}
\centering
\begin{tabular}{lcccc}
\toprule
\textbf{Method} & \textbf{Acc.} & \textbf{F1 (G/F/AI)} & \textbf{AUC} & \textbf{Time} \\
\midrule
\textbf{CertiVerify (Full)}      & \textbf{94.1 $\pm$ 0.8} & \textbf{0.95 / 0.91 / 0.89} & \textbf{0.95} & 2.6~s \\
ViT-Base (fine-tuned)            & 86.3 $\pm$ 1.1          & 0.89 / 0.84 / 0.73          & 0.88          & 1.1~s \\
GPT-4o (zero-shot)$^{*}$         & 69.3 $\pm$ 2.4          & 0.78 / 0.64 / 0.53          & 0.73          & 8.1~s \\
ELA + rule-based                 & 72.8 $\pm$ 1.7          & 0.81 / 0.67 / 0.58          & 0.76          & 0.9~s \\
Generic OCR + regex              & 76.4 $\pm$ 1.5          & 0.84 / 0.71 / 0.62          & 0.80          & 1.4~s \\
Text-only RoBERTa                & 80.2 $\pm$ 1.2          & 0.87 / 0.76 / 0.69          & 0.84          & 0.7~s \\
\bottomrule
\multicolumn{5}{l}{\footnotesize $^{*}$External API; not deployable under DPDP privacy constraints.}
\end{tabular}
\end{table}

The ViT-Base baseline is the critical comparison: fine-tuned on the same CertiForge-HE data, it isolates the effect of architectural choice from training-data advantage. Its plateau at 86.3\,\% and AI-Edited F1 of 0.73 (vs.\ CertiVerify's 0.89) confirm that a single holistic image classifier cannot jointly reason over pixel-level ELA anomalies, spatial layout deviation, and field-level semantic mismatch. Agent decomposition yields a 7.8-point gain attributable specifically to architectural separation of concerns.

\begin{table}[t]
\caption{Ablation study (AUC and class-specific F1).}
\label{tab:ablation}
\centering
\begin{tabular}{lccc}
\toprule
\textbf{Variant} & \textbf{AUC} & \textbf{Forged F1} & \textbf{AI-Edited F1} \\
\midrule
Full system                              & \textbf{0.95} & \textbf{0.91} & \textbf{0.89} \\
$-$ Visual Forensics Agent               & 0.82 & 0.74 & 0.61 \\
$-$ Field Cross-Check Agent              & 0.88 & 0.86 & 0.84 \\
$-$ Metadata Agent                       & 0.91 & 0.88 & 0.86 \\
$-$ Layout Agent                         & 0.90 & 0.87 & 0.83 \\
$-$ Reasoning Agent (avg.\ ensemble)     & 0.89 & 0.85 & 0.82 \\
Text Agent only                          & 0.84 & 0.76 & 0.69 \\
\bottomrule
\end{tabular}
\end{table}

Visual Forensics contributes the largest individual gain ($+0.13$ AUC); Field Cross-Check is second ($+0.07$ AUC) and is the only agent that detects genuine-certificate/incorrect-form-data mismatches. Replacing the learned Reasoning Agent with a simple average ensemble reduces AUC by 0.06.

\subsubsection{Visual vs.\ Layout Complementarity.} On the AI-Edited subset, Visual Forensics alone detects 71.1\,\%, Layout alone 58.6\,\%, both simultaneously 52.0\,\%, and the union 77.7\,\%. Only 52\,\% overlap confirms non-redundant signals: ELA captures frequency-domain artifacts from diffusion inpainting, while LayoutLMv3 captures spatial-domain character-spacing and line-height deviation. The full system reaches 89.0\,\% AI-Edited detection.

\subsubsection{Duplicate Detection.} The two-stage pipeline detects 97.8\,\% of duplicates with a 2.1\,\% false-positive rate, reducing false positives by 6.4 points over single-stage pgvector on the multi-day workshop subset.

\subsubsection{Workflow Impact and Explainability.} Verdict distribution on the full dataset: 64.7\,\% \texttt{auto\_verified}, 22.8\,\% \texttt{needs\_review}, 12.5\,\% \texttt{rejected}. Operational false-positive rate (genuine rejected) is 4.1\,\% [95\,\% CI: 3.2--5.1\,\%]. Average faculty review time drops from 4.2 to 1.2~minutes per submission (on the \texttt{needs\_review} subset), with an estimated 163 faculty-hours saved per 500-student cohort annually. Twelve faculty raters scored natural-language explanations for 50 \texttt{needs\_review} verdicts at a mean 4.6/5 on actionability (5-point Likert, $\sigma = 0.4$).

\section{Student Engagement Score}
\label{sec:ses}

Beyond verification, the platform computes a 0--100 Student Engagement Score (SES) from each student's verified portfolio. The base component weights categories (Internships and Leadership: 10; Academic Excellence: 9; Competitions: 8; Certifications: 7; Conferences: 6; Community Service: 5; Clubs: 4) with diminishing returns capped at three records per category. A diversity bonus adds 5 points per unique category (max 40). A verification-rate multiplier (0.2--1.0) scales the result by the fraction of submissions that are verified, penalising portfolios with predominantly unverified records. Grade bands: A+ ($\geq 90$), A ($\geq 75$), B+ ($\geq 60$), B ($\geq 45$), C ($\geq 30$), D ($< 30$).

\section{Security, Privacy, and Bias}
\label{sec:security}

Documents are stored in Supabase Storage with short-lived signed URLs; no certificate data leaves the institutional deployment. Row-level security enforces department-scoped faculty access. All agent outputs and verdicts are logged immutably for audit. The system is designed for compliance with the Indian DPDP Act 2023~\cite{dpdp2023}.

CertiVerify defends against naive forgery, AI-assisted forgery, and duplicate resubmission; it does not defend against insider collusion or adversarial perturbation attacks crafted against known weights (mitigated by keeping weights private to the institutional deployment). Two fairness risks were evaluated explicitly. \textit{Name and script bias:} NER \texttt{STUDENT\_NAME} F1 ranges from 0.93 (anglicised) to 0.89 (Northeast Indian and Dravidian-script transliterations)---a 4-point gap targeted by ongoing data collection. \textit{Issuer trust bias:} a naive prominence ranking would disadvantage regional and vernacular-event participants, so CertiVerify separates verification (does the certificate exist as claimed?) from prestige weighting, which is handled at the SES stage with explicit, auditable category weights.

\section{Future Work}
\label{sec:future}

Priority extensions: (i) NER coverage across all 22 scheduled Indian languages with native-script annotation; (ii) direct DigiLocker integration~\cite{digilocker2024} for near-machine-verifiable provenance on government-issued credentials; (iii) a continual-learning loop against GAN-generated adversarial forgeries; (iv) multi-institution federation for portable verified achievement records that survive institutional transfers.

\section{Conclusion}
\label{sec:conclusion}

CertiVerify decomposes certificate verification across five parallel specialist agents---text semantics, visual forensics, layout, metadata, and field-level cross-checking---synthesised by a calibrated Reasoning Agent. On CertiForge-HE it achieves 94.1\,\% accuracy and 0.95 AUC, outperforming a fine-tuned ViT-Base baseline by 7.8 points and zero-shot multimodal LLMs by 18--24 points. The framework automates verification for 64.7\,\% of submissions and provides explainable evidence for the remainder, reducing faculty workload by an estimated 71\,\% while improving record integrity and accreditation readiness. Domain-specific NER for Indian academic certificates, two-stage duplicate detection, non-overlapping forensic signals, and audit-ready explanations together establish a deployable foundation for trustworthy student activity management in Indian higher education.

\begin{thebibliography}{99}

\bibitem{balasubramanian2017}
Balasubramanian, S., Govindasamy, S.: Web-Based Student Information System for Managing Student Records. Int.\ J.\ Comput.\ Appl.\ \textbf{174}(2), 212--234 (2017)

\bibitem{alqahtani2020}
Alqahtani, A., Goodwin, S.: E-Portfolio Systems in Higher Education: Opportunities and Challenges. Educ.\ Inf.\ Technol.\ \textbf{25}(3), 1695--1712 (2020)

\bibitem{khandare2021}
Khandare, R., Jadhav, P.: Web-Based Activity Tracking System for Students. Int.\ J.\ Eng.\ Res.\ Technol.\ \textbf{10}(7), 568--573 (2021)

\bibitem{guillaro2023}
Guillaro, F., Cozzolino, D., Sud, A., Dufour, N., Verdoliva, L.: TruFor: Leveraging All-Round Information for Universal Image Forgery Detection and Localization. In: Proc.\ CVPR (2023)

\bibitem{kwon2022}
Kwon, M., Nam, S., Yu, I., Lee, H., Kim, C.: Learning JPEG Compression Artifacts for Image Manipulation Detection and Localization. Int.\ J.\ Comput.\ Vis.\ \textbf{130}, 1875--1895 (2022)

\bibitem{zhang2024}
Zhang, Y., Yu, S., Liu, J., Wu, X., et al.: ForgeryGPT: Multimodal Large Language Model for Explainable Image Forgery Detection and Localization. arXiv:2410.10238 (2024)

\bibitem{guo2025}
Guo, R., Tang, Y., Wang, M., et al.: M2F2-Det: Multi-Modal Interpretable Forged Face Detector. In: Proc.\ CVPR (2025)

\bibitem{chen2025}
Chen, X., Li, Z., Sun, P., et al.: Can Multi-modal LLMs Detect Document Manipulation? A Benchmark Study. arXiv:2508.11021 (2025)

\bibitem{docagent2025}
Chen, J., Wang, Y., Li, Z., et al.: DocAgent: Agentic Long Document Understanding with Multimodal LLMs. In: Proc.\ ACL Findings (2025)

\bibitem{spacy2017}
Honnibal, M., Montani, I.: spaCy: Industrial-Strength Natural Language Processing in Python. \url{https://spacy.io} (2017)

\bibitem{du2020}
Du, Y., Li, C., Guo, R., et al.: PP-OCR: A Practical Ultra Lightweight OCR System. arXiv:2009.09941 (2020)

\bibitem{breiman2001}
Breiman, L.: Random Forests. Mach.\ Learn.\ \textbf{45}(1), 5--32 (2001)

\bibitem{xu2022}
Xu, Y., Lv, T., Cui, L., et al.: LayoutLMv3: Pre-training for Document AI with Unified Text and Image Masking. In: Proc.\ ACM Multimedia (2022)

\bibitem{indic2020}
Kakwani, D., Kunchukuttan, A., Golla, S., et al.: IndicNLPSuite: Monolingual Corpora, Evaluation Benchmarks and Pre-trained Multilingual Language Models for Indian Languages. In: Findings of EMNLP (2020)

\bibitem{dosovitskiy2021}
Dosovitskiy, A., Beyer, L., Kolesnikov, A., et al.: An Image is Worth 16x16 Words: Transformers for Image Recognition at Scale. In: Proc.\ ICLR (2021)

\bibitem{he2016}
He, K., Zhang, X., Ren, S., Sun, J.: Deep Residual Learning for Image Recognition. In: Proc.\ CVPR, pp.\ 770--778 (2016)

\bibitem{selvaraju2017}
Selvaraju, R.R., Cogswell, M., Das, A., et al.: Grad-CAM: Visual Explanations from Deep Networks via Gradient-Based Localization. In: Proc.\ ICCV, pp.\ 618--626 (2017)

\bibitem{dpdp2023}
Ministry of Electronics and Information Technology, Government of India: Digital Personal Data Protection Act (2023)

\bibitem{naac2023}
National Assessment and Accreditation Council: Manual for Self-Study Report---Affiliated/Constituent Colleges. NAAC, Bangalore (2023)

\bibitem{mitchell2023}
Mitchell, E., Lee, Y., Khazatsky, A., et al.: DetectGPT: Zero-Shot Machine-Generated Text Detection using Probability Curvature. In: Proc.\ ICML (2023)

\bibitem{cozzolino2020}
Cozzolino, D., Verdoliva, L.: Noiseprint: A CNN-Based Camera Model Fingerprint. IEEE Trans.\ Inf.\ Forensics Secur.\ \textbf{15}, 144--159 (2020)

\bibitem{digilocker2024}
DigiLocker, Ministry of Electronics and Information Technology: DigiLocker Issuer API Specification v3.2. Government of India (2024)

\end{thebibliography}

\end{document}
